%% ================================================================================================
%%  Appendix B. Algorithms
%%  Code: simulations/python/voting.py (schulze, split_cycle, ranked_pairs, qvc, pvc_member)
%% ================================================================================================
\section{Algorithms}\label{app:algorithms}

The reference implementation is \texttt{simulations/python/voting.py}; its \texttt{--check} option tests
each algorithm below against a brute-force or textbook result (Appendix~\ref{app:simulation}).

\paragraph{Pairwise counts.} $\Maj(x,y) = |\{i : x \pref_i y\}|$ for all pairs, in $O(nm^2)$ operations.
Ties in a ballot count for neither option.

\paragraph{Schulze method} \citep{Schulze2011}. On a set $S$ of options, the link from $x$ to $y$ has
strength $d(x,y) = \Maj(x,y)$ if $\Maj(x,y) > \Maj(y,x)$ and $0$ otherwise (winning votes). The strength
$p(x,y)$ of the strongest path from $x$ to $y$ is the largest, over paths $x = c_0, c_1, \dots, c_k = y$ in
$S$, of the weakest link $\min_j d(c_{j-1}, c_j)$. It is computed by the Floyd--Warshall recursion for
widest paths:
\begin{equation*}
p(x,y) \gets \max\bigl(p(x,y),\ \min(p(x,z), p(z,y))\bigr) \quad \text{for } z \in S, \text{ then all } x, y,
\end{equation*}
starting from $p = d$, in $O(|S|^3)$ operations. The winners are the options $x$ with
$p(x,y) \ge p(y,x)$ for every $y \in S$; there is always at least one.

\paragraph{Split Cycle} \citep{HollidayPacuit2023splitcycle}. With margins $M(x,y) = \Maj(x,y) -
\Maj(y,x)$, let $q(y,x)$ be the strength of the strongest path from $y$ to $x$ through positive margins,
computed as above. Option $x$ defeats $y$ if $M(x,y) > 0$ and $M(x,y) > q(y,x)$: the majority of $x$ over
$y$ is not the weakest link of a cycle. The winners are the undefeated options.

\paragraph{Ranked Pairs} \citep{Tideman1987}. Sort the pairwise majorities by decreasing strength and lock
each one in turn unless it would close a cycle with those already locked; the winner is the option with no
locked defeat.

\paragraph{\QVC($\theta$).} Algorithm~\ref{alg:qvc}: veto counts in $O(nm)$, the admissible set in $O(m)$,
the Schulze winners of the admissible set in $O(nm^2 + m^3)$.

\paragraph{Proportional veto core.} Testing every coalition is exponential in $n$. A max-flow test is
polynomial; \citet{IanovskiKondratev2021} give a polynomial algorithm, and the test below is the one we
implement.

\begin{lemma}\label{lem:pvc-flow}
For option $x$ and voter $i$, let $W_i(x)$ be the set of options that $i$ ranks at or below $x$ ($x$
included). Then $x$ is in the proportional veto core if and only if, for every coalition $T$,
$\bigl|\bigcup_{i \in T} W_i(x)\bigr| \ge m |T| / n$. Equivalently, in the network with a source, one
node per voter and per option, and a sink, with arcs source $\to i$ of capacity $m$, $i \to y$ of
unbounded capacity for $y \in W_i(x)$, and $y \to$ sink of capacity $n$, the maximum flow equals $mn$.
\end{lemma}

\begin{proof}
Coalition $T$ of $t$ voters blocks $x$ if at least $m - \lceil mt/n \rceil + 1$ options are preferred to
$x$ by all its members. The options preferred to $x$ by every member of $T$ are the complement of
$\bigcup_{i \in T} W_i(x)$, so $T$ blocks $x$ if and only if
$m - |\bigcup_T W_i(x)| \ge m - \lceil mt/n \rceil + 1$, that is
$|\bigcup_T W_i(x)| < \lceil mt/n \rceil$, which for an integer left-hand side is
$|\bigcup_T W_i(x)| < mt/n$. For the flow, a cut that keeps the voters of $T$ on the source side must also
keep $\bigcup_T W_i(x)$ there (the arcs from voters are unbounded); its capacity is
$m(n-t) + n\,|\bigcup_T W_i(x)|$. The minimum cut, hence the maximum flow, equals $mn$ if and only if
$n\,|\bigcup_T W_i(x)| \ge mt$ for every $T$. \qedhere
\end{proof}

The test takes one max-flow computation per option, on a network of $n + m + 2$ nodes. Our check compares
it with the definition, by enumeration of all coalitions, on random profiles with up to seven voters and
six options, with and without ties.

\paragraph{Tie-breaking.} Every rule in the simulations returns a set of tied winners, broken by fewest
vetoes, then the Schulze winners among the tied options, then the lowest index (the order of options is
random). Both criteria are independent of clones, so the tie-break does not undo that property where a
rule has it.
